% ============================================================================
\section{Length Normalization Changes the Training Mixture}
\label{sec:normalization}
% ============================================================================

Length imbalance depends on how token losses are reduced. MOPD writes a
response-normalized objective
\citep{ma2026mopdmultiteacheronpolicydistillation}, whereas Open-MOPD
analyzes a token-mean implementation in which instruction following supplies
approximately $20\%$ of prompts but only $1\%$ of response tokens
\citep{gao2026openmopddiagnosingfixingcapability}. This observation motivates
separating the training objective from the sampling mixture. It does not
establish a length pathology intrinsic to multi-teacher distillation.
The formal objective and each implementation's complete reduction path,
including gradient accumulation, require separate verification.
\missing{Reduction audit.}

\paragraph{Three reduction orders.}
Fix a rollout batch containing $n_d\geq1$ responses from domain $d$.
Response $i$ has $T_{di}>0$ valid positions, position losses $\ell_{dit}$,
and mean loss $\bar\ell_{di}=T_{di}^{-1}\sum_t\ell_{dit}$.
Write $N_d=\sum_iT_{di}$, $N=\sum_dN_d$, and $s_d=N_d/N$.
The target domain weights $a_d\geq0$, with $\sum_da_d=1$, are distinct
from the prompt proportions $p_d=n_d/\sum_jn_j$. Consider
\begin{align}
 L_{\mathrm{tok}}
   &=\frac1N\sum_{d,i,t}\ell_{dit}
     =\sum_ds_d\frac{\sum_iT_{di}\bar\ell_{di}}{N_d},
     \label{eq:norm_global}\\
 L_{\mathrm{DT}}
   &=\sum_da_d\frac{\sum_iT_{di}\bar\ell_{di}}{N_d},
     \label{eq:norm_domain_token}\\
 L_{\mathrm{DR}}
   &=\sum_da_d\frac1{n_d}\sum_i\bar\ell_{di}.
     \label{eq:norm_domain_response}
\end{align}
These respectively average tokens globally, tokens within domains, and
positions within responses before averaging responses within domains.
Multiplying each domain's token loss by $a_d/s_d$ converts
Equation~\ref{eq:norm_global} into Equation~\ref{eq:norm_domain_token};
this is Open-MOPD's token-share balancing. It preserves length weighting
within each domain. An ordinary mean over response means instead equals
Equation~\ref{eq:norm_domain_response} with $a_d=p_d$.

\begin{proposition}[Fixed-batch length--gradient decomposition]
\label{prop:norm_covariance}
Hold the rollouts, masks, and domain weights fixed when differentiating;
any surrogate advantages are detached. Define
$h_{di}=T_{di}^{-1}\sum_t\nabla_\theta\ell_{dit}$,
$\bar h_d=n_d^{-1}\sum_ih_{di}$, and
$\bar T_d=n_d^{-1}\sum_iT_{di}$. Let
\begin{equation}
 C_d=\frac1{n_d}\sum_i(T_{di}-\bar T_d)(h_{di}-\bar h_d).
 \label{eq:norm_covariance}
\end{equation}
Then
\begin{align}
 \nabla L_{\mathrm{tok}}
   &=\sum_ds_d\left(\bar h_d+\frac{C_d}{\bar T_d}\right),
     \label{eq:norm_global_gradient}\\
 \nabla L_{\mathrm{DT}}
   &=\sum_da_d\left(\bar h_d+\frac{C_d}{\bar T_d}\right),
 &\nabla L_{\mathrm{DR}}&=\sum_da_d\bar h_d.
     \label{eq:norm_domain_gradients}
\end{align}
\end{proposition}
\begin{proof}
Expanding the empirical covariance gives
$n_d^{-1}\sum_iT_{di}h_{di}=\bar T_d\bar h_d+C_d$.
Divide by $\bar T_d$ and substitute into each reduction.
\end{proof}

Domain token normalization therefore removes the coefficient mismatch
$s_d-a_d$ while retaining the within-domain length--gradient covariance.
These are standard weighted-mean identities, used here to explain existing
training recipes rather than introduce a new algorithm. Equal token
counts do not imply equal gradient norms or compatible update directions;
neither token shares nor scalar reward magnitudes capture vector
cancellation or optimizer preconditioning. The identities are exact for
a fixed batch. In expectation, a sample ratio generally differs from
$\mathbb E[Th]/\mathbb E[T]$; the latter describes an appropriate
large-batch limit, not an exact finite-batch expectation.

\paragraph{Controlled observations.}
To test this explanation, compare all three reductions on identical
rollouts, teacher scores, and surrogate signals at early, middle, and late
checkpoints. Measure $s_d$, $\bar h_d$, $C_d/\bar T_d$, and gradient
directions before clipping. First set $a_d=p_d$, or use equal prompt
quotas, to isolate length weighting; evaluate uniform target weights
separately. Otherwise, correcting length also changes the domain prior.
\missing{Gradient decomposition.}
Follow these measurements with matched training branches that retain the
same prompt schedule, response limits, teachers, and update settings,
while generating their own on-policy responses. Report capability curves
against both updates and cumulative generated tokens.
\missing{Training curves.}
Reducing the response limit is a separate intervention: it changes the
available prefixes and may remove solution endings. Its effect cannot be
identified with loss reweighting, nor do the identities imply that
response normalization always improves capability integration.
